\section{Homogeneous Linear Equations}

\noindent
A second order differential equation is said to be {\color{blue}\it linear\/} if it can be written
as
\begin{equation}
  \label{eq:5.1.1}y''+p(x)y'+q(x)y=f(x).
\end{equation}
We call the function $f$ on the right a {\color{blue}\it forcing function\/}, since in physical
applications it's often related to a force acting on some system modeled by the differential
equation. We say that \eqref{eq:5.1.1} is {\color{blue}\it homogeneous\/} if $f\equiv0$ or
{\color{blue}\it nonhomogeneous\/} if $f\not\equiv0$. Since these definitions are like the
corresponding definitions in Section~2.1 for the linear first order equation
\begin{equation}
  \label{eq:5.1.2}y'+p(x)y=f(x),
\end{equation}
it's natural to expect similarities between methods of solving \eqref{eq:5.1.1} and
\eqref{eq:5.1.2}. However, solving \eqref{eq:5.1.1} is more difficult than solving \eqref{eq:5.1.2}.
For example, while Theorem~\ref{thmtype:2.1.1} gives a formula for the general solution of
\eqref{eq:5.1.2} in the case where $f\equiv0$ and Theorem~\ref{thmtype:2.1.2} gives a formula for
the case where $f\not\equiv0$, there are no formulas for the general solution of \eqref{eq:5.1.1} in
either case. Therefore we must be content to solve linear second order equations of special forms.

In Section~2.1 we considered the homogeneous equation $y'+p(x)y=0$ first, and then used a nontrivial
solution of this equation to find the general solution of the nonhomogeneous equation
$y'+p(x)y=f(x)$. Although the progression from the homogeneous to the nonhomogeneous case isn't that
simple for the linear second order equation, it's still necessary to solve the homogeneous equation
\begin{equation}
  \label{eq:5.1.3}y''+p(x)y'+q(x)y=0
\end{equation}
in order to solve the nonhomogeneous equation \eqref{eq:5.1.1}. This section is devoted to
\eqref{eq:5.1.3}.

The next theorem gives sufficient conditions for existence and uniqueness of solutions of initial
value problems for \eqref{eq:5.1.3}. We omit the proof.

\begin{theorem}
  \color{blue}
  \label{thmtype:5.1.1} Suppose $p$ and $q$ are continuous on an open interval $(a,b)$, let $x_{0}$
  be any point in $(a,b)$, and let $k_{0}$ and $k_{1}$ be arbitrary real numbers. Then the initial
  value problem
  \[
    y''+p(x)y'+q(x)y=0,\ y(x_{0})=k_{0},\ y'(x_{0})=k_{1}
  \]
  has a unique solution on $(a,b)$.
\end{theorem}

Since $y\equiv0$ is obviously a solution of \eqref{eq:5.1.3} we call it the {\color{blue}\it
trivial} solution. Any other solution is {\color{blue}\it nontrivial}. Under the assumptions of
Theorem~\ref{thmtype:5.1.1}, the only solution of the initial value problem
\[
  y''+p(x)y'+q(x)y=0,\ y(x_{0})=0,\ y'(x_{0})=0
\]
on $(a,b)$ is the trivial solution (Exercise~\ref{exer:5.1.24}).

The next three examples illustrate concepts that we'll develop later in this section. You shouldn't
be concerned with how to {\color{blue}\it find\/} the given solutions of the equations in these
examples. This will be explained in later sections.

\begin{example}
  \label{example:5.1.1} \rm The coefficients of $y'$ and $y$ in
  \begin{equation}
    \label{eq:5.1.4}y''-y=0
  \end{equation}
  are the constant functions $p\equiv0$ and $q\equiv-1$, which are continuous on $(-\infty,\infty)$.
  Therefore Theorem~\ref{thmtype:5.1.1} implies that every initial value problem for
  \eqref{eq:5.1.4} has a unique solution on $(-\infty,\infty)$.
  \begin{alist}
    \item % (a)
    Verify that $y_{1}=e^{x}$ and $y_{2}=e^{-x}$ are solutions of \eqref{eq:5.1.4} on
    $(-\infty,\infty)$.

    \item % (b)
    Verify that if $c_{1}$ and $c_{2}$ are arbitrary constants, $y=c_{1}e^{x}+c_{2}e^{-x}$ is a
    solution of \eqref{eq:5.1.4} on $(-\infty,\infty)$.

    \item % (c)
    Solve the initial value problem
    \begin{equation}
      \label{eq:5.1.5}y''-y=0,\quad y(0)=1,\quad y'(0)=3.
    \end{equation}
  \end{alist}
\end{example}

\solutionpart{a} If $y_{1}=e^{x}$ then $y_{1}'=e^{x}$ and $y_{1}''=e^{x}=y_{1}$, so
$y_{1}''-y_{1}=0$. If $y_{2}=e^{-x}$, then $y_{2}'=-e^{-x}$ and $y_{2}''=e^{-x}=y_{2}$, so
$y_{2}''-y_{2}=0$.

\solutionpart{b} If
\begin{equation}
  \label{eq:5.1.6}y=c_{1}e^{x}+c_{2}e^{-x}
\end{equation}
then
\begin{equation}
  \label{eq:5.1.7}y'=c_{1}e^{x}-c_{2}e^{-x}
\end{equation}
and
\[
  y''=c_{1}e^{x}+c_{2}e^{-x},
\]
so
\begin{eqnarray*}
  y''-y & = & (c_{1}e^{x}+c_{2}e^{-x})-(c_{1}e^{x}+c_{2}e^{-x}) \\
        & = & c_{1}(e^{x}-e^{x})+c_{2}(e^{-x}-e^{-x})=0
\end{eqnarray*}
for all $x$. Therefore $y=c_{1}e^{x}+c_{2}e^{-x}$ is a solution of \eqref{eq:5.1.4} on
$(-\infty,\infty)$.

\solutionpart{c} We can solve \eqref{eq:5.1.5} by choosing $c_{1}$ and $c_{2}$ in \eqref{eq:5.1.6}
so that $y(0)=1$ and $y'(0)=3$. Setting $x=0$ in \eqref{eq:5.1.6} and \eqref{eq:5.1.7} shows that
this is equivalent to
\begin{eqnarray*}
  c_{1}+c_{2} & = & 1 \\
  c_{1}-c_{2} & = & 3.
\end{eqnarray*}
Solving these equations yields $c_{1}=2$ and $c_{2}=-1$. Therefore $y=2e^{x}-e^{-x}$ is the unique
solution of \eqref{eq:5.1.5} on $(-\infty,\infty)$.

\begin{example}
  \label{example:5.1.2} \rm %\newcommand{\exampletrig}
  Let $\omega$ be a positive constant. The coefficients of $y'$ and $y$ in
  \begin{equation}
    \label{eq:5.1.8}y''+\omega^{2}y=0
  \end{equation}
  are the constant functions $p\equiv0$ and $q\equiv\omega^{2}$, which are continuous on
  $(-\infty,\infty)$. Therefore Theorem~\ref{thmtype:5.1.1} implies that every initial value problem
  for \eqref{eq:5.1.8} has a unique solution on $(-\infty,\infty)$.
  \begin{alist}
    \item % (a)
    Verify that $y_{1}=\cos\omega x$ and $y_{2}=\sin\omega x$ are solutions of \eqref{eq:5.1.8} on
    $(-\infty,\infty)$.

    \item % (b)
    Verify that if $c_{1}$ and $c_{2}$ are arbitrary constants then
    $y=c_{1}\cos\omega x+c_{2}\sin\omega x$ is a solution of \eqref{eq:5.1.8} on $(-\infty,\infty)$.

    \item % (c)
    Solve the initial value problem
    \begin{equation}
      \label{eq:5.1.9}y''+\omega^{2}y=0,\quad y(0)=1,\quad y'(0)=3.
    \end{equation}
  \end{alist}
\end{example}

\solutionpart{a} If $y_{1}=\cos\omega x$ then $y_{1}'=-\omega\sin\omega x$ and
$y_{1}''=-\omega^{2}\cos\omega x=-\omega^{2}y_{1}$, so $y_{1}''+\omega^{2}y_{1}=0$. If
$y_{2}=\sin\omega x$ then, $y_{2}'=\omega\cos\omega x$ and
$y_{2}''=-\omega^{2}\sin\omega x=-\omega^{2}y_{2}$, so $y_{2}''+\omega^{2}y_{2}=0$.

\solutionpart{b} If
\begin{equation}
  \label{eq:5.1.10}y=c_{1}\cos\omega x+c_{2}\sin\omega x
\end{equation}
then
\begin{equation}
  \label{eq:5.1.11}y'=\omega(-c_{1}\sin\omega x+c_{2}\cos\omega x)
\end{equation}
and
\[
  y''=-\omega^{2}(c_{1}\cos\omega x+c_{2}\sin\omega x),
\]
so
\begin{eqnarray*}
  y''+\omega^{2}y & = & -\omega^{2}(c_{1}\cos\omega x+c_{2}\sin\omega x)+\omega^{2}(c_{1}\cos\omega x+c_{2}\sin\omega x) \\
                  & = & c_{1}\omega^{2}(-\cos\omega x+\cos\omega x)+c_{2}\omega^{2}(-\sin\omega x+\sin\omega x)=0
\end{eqnarray*}
for all $x$. Therefore $y=c_{1}\cos\omega x+c_{2}\sin\omega x$ is a solution of \eqref{eq:5.1.8} on
$(-\infty,\infty)$.

\solutionpart{c} To solve \eqref{eq:5.1.9}, we must choosing $c_{1}$ and $c_{2}$ in
\eqref{eq:5.1.10} so that $y(0)=1$ and $y'(0)=3$. Setting $x=0$ in \eqref{eq:5.1.10} and
\eqref{eq:5.1.11} shows that $c_{1}=1$ and $c_{2}=3/\omega$. Therefore
\[
  y=\cos\omega x+\frac{3}{\omega}\sin\omega x
\]
is the unique solution of \eqref{eq:5.1.9} on $(-\infty,\infty)$. \bbox

Theorem~\ref{thmtype:5.1.1} implies that if $k_{0}$ and $k_{1}$ are arbitrary real numbers then the
initial value problem
\begin{equation}
  \label{eq:5.1.12}P_{0}(x)y''+P_{1}(x)y'+P_{2}(x)y=0,\quad y(x_{0})=k_{0},\quad y'(x_{0})=k_{1}
\end{equation}
has a unique solution on an interval $(a,b)$ that contains $x_{0}$, provided that $P_{0}$, $P_{1}$,
and $P_{2}$ are continuous and $P_{0}$ has no zeros on $(a,b)$. To see this, we rewrite the
differential equation in \eqref{eq:5.1.12} as
\[
  y''+\frac{P_{1}(x)}{P_{0}(x)}y'+\frac{P_{2}(x)}{P_{0}(x)}y=0
\]
and apply Theorem~\ref{thmtype:5.1.1} with $p=P_{1}/P_{0}$ and $q=P_{2}/P_{0}$.

\begin{example}
  \label{example:5.1.3} \rm The equation
  \begin{equation}
    \label{eq:5.1.13}x^{2}y''+xy'-4y=0
  \end{equation}
  has the form of the differential equation in \eqref{eq:5.1.12}, with $P_{0}(x)=x^{2}$,
  $P_{1}(x)=x$, and $P_{2}(x)=-4$, which are all continuous on $(-\infty,\infty)$. However, since
  $P(0)=0$ we must consider solutions of \eqref{eq:5.1.13} on $(-\infty,0)$ and $(0,\infty)$. Since
  $P_{0}$ has no zeros on these intervals, Theorem~\ref{thmtype:5.1.1} implies that the initial
  value problem
  \[
    x^{2}y''+xy'-4y=0,\quad y(x_{0})=k_{0},\quad y'(x_{0})=k_{1}
  \]
  has a unique solution on $(0,\infty)$ if $x_{0}>0$, or on $(-\infty,0)$ if $x_{0}<0$.
  \begin{alist}
    \item % (a)
    Verify that $y_{1}=x^{2}$ is a solution of \eqref{eq:5.1.13} on $(-\infty,\infty)$ and
    $y_{2}=1/x^{2}$ is a solution of \eqref{eq:5.1.13} on $(-\infty,0)$ and $(0,\infty)$.

    \item % (b)
    Verify that if $c_{1}$ and $c_{2}$ are any constants then $y=c_{1}x^{2}+c_{2}/x^{2}$ is a
    solution of \eqref{eq:5.1.13} on $(-\infty,0)$ and $(0,\infty)$.

    \item % (c)
    Solve the initial value problem
    \begin{equation}
      \label{eq:5.1.14}x^{2}y''+xy'-4y=0,\quad y(1)=2,\quad y'(1)=0.
    \end{equation}

    \item % (d)
    Solve the initial value problem
    \begin{equation}
      \label{eq:5.1.15}x^{2}y''+xy'-4y=0,\quad y(-1)=2,\quad y'(-1)=0.
    \end{equation}
  \end{alist}
\end{example}

\solutionpart{a} If $y_{1}=x^{2}$ then $y_{1}'=2x$ and $y_{1}''=2$, so
\[
  x^{2}y_{1}''+xy_{1}'-4y_{1}=x^{2}(2)+x(2x)-4x^{2}=0
\]
for $x$ in $(-\infty,\infty)$. If $y_{2}=1/x^{2}$, then $y_{2}'=-2/x^{3}$ and $y_{2}''=6/x^{4}$, so
\[
  x^{2}y_{2}''+xy_{2}'-4y_{2}=x^{2}\left(\frac{6}{x^{4}}\right)-x\left(\frac{2}{x^{3}}\right)-\frac{4}{x^{2}}=0
\]
for $x$ in $(-\infty,0)$ or $(0,\infty)$.

\solutionpart{b} If
\begin{equation}
  \label{eq:5.1.16}y=c_{1}x^{2}+\frac{c_{2}}{x^{2}}
\end{equation}
then
\begin{equation}
  \label{eq:5.1.17}y'=2c_{1}x-\frac{2c_{2}}{x^{3}}
\end{equation}
and
\[
  y''=2c_{1}+\frac{6c_{2}}{x^{4}},
\]
so
\begin{eqnarray*}
  x^{2}y''+xy'-4y & = & x^{2}\dst{\left(2c_{1}+\frac{6c_{2}}{x^{4}}\right)}+x\dst{\left(2c_{1}x-\frac{2c_{2}}{x^{3}}\right)}-4\dst{\left(c_{1}x^{2}+\frac{c_{2}}{x^{2}}\right)} \\[2\jot]
                  & = & c_{1}(2x^{2}+2x^{2}-4x^{2})+c_{2}\dst{\left(\frac{6}{x^{2}}-\frac{2}{x^{2}}-\frac{4}{x^{2}}\right)}                                                     \\[2\jot]
                  & = & c_{1}\cdot0+c_{2}\cdot0=0
\end{eqnarray*}
for $x$ in $(-\infty,0)$ or $(0,\infty)$.

\solutionpart{c} To solve \eqref{eq:5.1.14}, we choose $c_{1}$ and $c_{2}$ in \eqref{eq:5.1.16} so
that $y(1)=2$ and $y'(1)=0$. Setting $x=1$ in \eqref{eq:5.1.16} and \eqref{eq:5.1.17} shows that
this is equivalent to
\begin{eqnarray*}
  \phantom{2}c_{1}+\phantom{2}c_{2} & = & 2 \\
  2c_{1}-2c_{2}                     & = & 0.
\end{eqnarray*}
Solving these equations yields $c_{1}=1$ and $c_{2}=1$. Therefore $y=x^{2}+1/x^{2}$ is the unique
solution of \eqref{eq:5.1.14} on $(0,\infty)$.

\solutionpart{d} We can solve \eqref{eq:5.1.15} by choosing $c_{1}$ and $c_{2}$ in \eqref{eq:5.1.16}
so that $y(-1)=2$ and $y'(-1)=0$. Setting $x=-1$ in \eqref{eq:5.1.16} and \eqref{eq:5.1.17} shows
that this is equivalent to
\begin{eqnarray*}
  \phantom{-2}c_{1}+\phantom{2}c_{2} & = & 2 \\
  -2c_{1}+2c_{2}                     & = & 0.
\end{eqnarray*}
Solving these equations yields $c_{1}=1$ and $c_{2}=1$. Therefore $y=x^{2}+1/x^{2}$ is the unique
solution of \eqref{eq:5.1.15} on $(-\infty,0)$. \bbox

Although the {\color{blue}\it formulas\/} for the solutions of \eqref{eq:5.1.14} and
\eqref{eq:5.1.15} are both $y=x^{2}+1/x^{2}$, you should not conclude that these two initial value
problems have the same solution. Remember that a solution of an initial value problem is defined
{\color{blue}\it on an interval that contains the initial point\/}; therefore, the solution of
\eqref{eq:5.1.14} is $y=x^{2}+1/x^{2}$ {\color{blue}\it on the interval\/} $(0,\infty)$, which
contains the initial point $x_{0}=1$, while the solution of \eqref{eq:5.1.15} is $y=x^{2}+1/x^{2}$
{\color{blue}\it on the interval\/} $(-\infty,0)$, which contains the initial point $x_{0}=-1$.

\boxit{The General Solution of a Homogeneous Linear Second Order Equation}

\noindent
If $y_{1}$ and $y_{2}$ are defined on an interval $(a,b)$ and $c_{1}$ and $c_{2}$ are constants,
then
\[
  y=c_{1}y_{1}+c_{2}y_{2}
\]
is a {\color{blue}\it linear combination of $y_{1}$ and $y_{2}$\/}. For example, $y=2\cos x+7\sin x$
is a linear combination of $y_{1}= \cos x$ and $y_{2}=\sin x$, with $c_{1}=2$ and $c_{2}=7$.

The next theorem states a fact that we've already verified in Examples~\ref{example:5.1.1},
\ref{example:5.1.2}, and \ref{example:5.1.3}.

\begin{theorem}
  \color{blue}
  \label{thmtype:5.1.2} If $y_{1}$ and $y_{2}$ are solutions of the homogeneous equation
  \begin{equation}
    \label{eq:5.1.18}y''+p(x)y'+q(x)y=0
  \end{equation}
  on $(a,b)$, then any linear combination
  \begin{equation}
    \label{eq:5.1.19}y=c_{1}y_{1}+c_{2}y_{2}
  \end{equation}
  of $y_{1}$ and $y_{2}$ is also a solution of \eqref{eq:5.1.18} on $(a,b)$.
\end{theorem}

\proof If
\[
  y=c_{1}y_{1}+c_{2}y_{2}
\]
then
\[
  y'=c_{1}y_{1}'+c_{2}y_{2}'\mbox{\quad and \quad}y''=c_{1}y_{1}''+c_{2}y_{2}''.
\]
Therefore
\begin{eqnarray*}
  y''+p(x)y'+q(x)y & = & (c_{1}y_{1}''+c_{2}y_{2}'')+p(x)(c_{1}y_{1}'+c_{2}y_{2}')+q(x)(c_{1}y_{1}+c_{2}y_{2})         \\[2\jot]
                   & = & c_{1}\left(y_{1}''+p(x)y_{1}'+q(x)y_{1}\right)+c_{2}\left(y_{2}''+p(x)y_{2}'+q(x)y_{2}\right) \\[2\jot]
                   & = & c_{1}\cdot0+c_{2}\cdot0=0,
\end{eqnarray*}
since $y_{1}$ and $y_{2}$ are solutions of \eqref{eq:5.1.18}. \bbox

We say that $\{y_{1},y_{2}\}$ is a {\color{blue}\it fundamental set of solutions of
\eqref{eq:5.1.18} on \/} $(a,b)$ if every solution of \eqref{eq:5.1.18} on $(a,b)$ can be written
as a linear combination of $y_{1}$ and $y_{2}$ as in \eqref{eq:5.1.19}. In this case we say that
\eqref{eq:5.1.19} is {\color{blue}\it general solution of \eqref{eq:5.1.18} on\/} $(a,b)$.

\boxit{Linear Independence}

\noindent
We need a way to determine whether a given set $\{y_{1},y_{2}\}$ of solutions of \eqref{eq:5.1.18}
is a fundamental set. The next definition will enable us to state necessary and sufficient
conditions for this.

We say that two functions $y_{1}$ and $y_{2}$ defined on an interval $(a,b)$ are {\color{blue}\it
linearly independent on\/} $(a,b)$ if neither is a constant multiple of the other on $(a,b)$. (In
particular, this means that neither can be the trivial solution of \eqref{eq:5.1.18}, since, for
example, if $y_{1}\equiv0$ we could write $y_{1}=0y_{2}$.) We'll also say that the set
$\{y_{1},y_{2}\}$ {\color{blue}\it is linearly independent on\/} $(a,b)$.

\begin{theorem}
  \color{blue}
  \label{thmtype:5.1.3} Suppose $p$ and $q$ are continuous on $(a,b)$. Then a set $\{y_{1},y_{2}\}$
  of solutions of
  \begin{equation}
    \label{eq:5.1.20}y''+p(x)y'+q(x)y=0
  \end{equation}
  on $(a,b)$ is a fundamental set if and only if $\{y_{1},y_{2}\}$ is linearly independent on
  $(a,b)$.
\end{theorem}

We'll present the proof of Theorem~\ref{thmtype:5.1.3} in steps worth regarding as theorems in their
own right. However, let's first interpret Theorem~\ref{thmtype:5.1.3} in terms of
Examples~\ref{example:5.1.1}, \ref{example:5.1.2}, and \ref{example:5.1.3}.

\begin{example}
  \label{example:5.1.4} \rm \mbox{}
  \newline
  \begin{alist}
    \item % (a)
    Since $e^{x}/e^{-x}=e^{2x}$ is nonconstant, Theorem~\ref{thmtype:5.1.3} implies that
    $y=c_{1}e^{x}+c_{2}e^{-x}$ is the general solution of $y''-y=0$ on $(-\infty,\infty)$.

    \item % (b)
    Since $\cos\omega x/\sin\omega x=\cot\omega x$ is nonconstant, Theorem~\ref{thmtype:5.1.3}
    implies that $y=c_{1}\cos\omega x+c_{2}\sin\omega x$ is the general solution of
    $y''+\omega^{2}y=0$ on $(-\infty,\infty)$.

    \item % (c)
    Since $x^{2}/x^{-2}=x^{4}$ is nonconstant, Theorem~\ref{thmtype:5.1.3} implies that
    $y=c_{1}x^{2}+c_{2}/x^{2}$ is the general solution of $x^{2}y''+xy'-4y=0$ on $(-\infty,0)$ and
    $(0,\infty)$.
  \end{alist}
\end{example}

\boxit{The Wronskian and Abel's Formula}

\noindent
To motivate a result that we need in order to prove Theorem~\ref{thmtype:5.1.3}, let's see what is
required to prove that $\{y_{1},y_{2}\}$ is a fundamental set of solutions of \eqref{eq:5.1.20} on
$(a,b)$. Let $x_{0}$ be an arbitrary point in $(a,b)$, and suppose $y$ is an arbitrary solution of
\eqref{eq:5.1.20} on $(a,b)$. Then $y$ is the unique solution of the initial value problem
\begin{equation}
  \label{eq:5.1.21}y''+p(x)y'+q(x)y=0,\quad y(x_{0})=k_{0},\quad y'(x_{0})=k_{1};
\end{equation}
that is, $k_{0}$ and $k_{1}$ are the numbers obtained by evaluating $y$ and $y'$ at $x_{0}$.
Moreover, $k_{0}$ and $k_{1}$ can be any real numbers, since Theorem~\ref{thmtype:5.1.1} implies
that \eqref{eq:5.1.21} has a solution no matter how $k_{0}$ and $k_{1}$ are chosen. Therefore
$\{y_{1},y_{2}\}$ is a fundamental set of solutions of \eqref{eq:5.1.20} on $(a,b)$ if and only if
it's possible to write the solution of an arbitrary initial value problem \eqref{eq:5.1.21} as
$y=c_{1}y_{1}+c_{2}y_{2}$. This is equivalent to requiring that the system
\begin{equation}
  \label{eq:5.1.22}
  \begin{array}{rcl}
    c_{1}y_{1}(x_{0})+c_{2}y_{2}(x_{0})   & = & k_{0} \\
    c_{1}y_{1}'(x_{0})+c_{2}y_{2}'(x_{0}) & = & k_{1}
  \end{array}
\end{equation}
has a solution $(c_{1},c_{2})$ for every choice of $(k_{0},k_{1})$. Let's try to solve
\eqref{eq:5.1.22}.

Multiplying the first equation in \eqref{eq:5.1.22} by $y_{2}'(x_{0})$ and the second by
$y_{2}(x_{0})$ yields
\begin{eqnarray*}
  c_{1}y_{1}(x_{0})y_{2}'(x_{0})+c_{2}y_{2}(x_{0})y_{2}'(x_{0}) & = & y_{2}'(x_{0})k_{0} \\
  c_{1}y_{1}'(x_{0})y_{2}(x_{0})+c_{2}y_{2}'(x_{0})y_{2}(x_{0}) & = & y_{2}(x_{0})k_{1},
\end{eqnarray*}
and subtracting the second equation here from the first yields
\begin{equation}
  \label{eq:5.1.23}\left(y_{1}(x_{0})y_{2}'(x_{0})-y_{1}'(x_{0})y_{2}(x_{0})\right)c_{1}= y_{2}'(x_{0})k_{0}-y_{2}(x_{0})k_{1}.
\end{equation}
Multiplying the first equation in \eqref{eq:5.1.22} by $y_{1}'(x_{0})$ and the second by
$y_{1}(x_{0})$ yields
\begin{eqnarray*}
  c_{1}y_{1}(x_{0})y_{1}'(x_{0})+c_{2}y_{2}(x_{0})y_{1}'(x_{0}) & = & y_{1}'(x_{0})k_{0} \\
  c_{1}y_{1}'(x_{0})y_{1}(x_{0})+c_{2}y_{2}'(x_{0})y_{1}(x_{0}) & = & y_{1}(x_{0})k_{1},
\end{eqnarray*}
and subtracting the first equation here from the second yields
\begin{equation}
  \label{eq:5.1.24}\left(y_{1}(x_{0})y_{2}'(x_{0})-y_{1}'(x_{0})y_{2}(x_{0})\right)c_{2}= y_{1}(x_{0})k_{1}-y_{1}'(x_{0})k_{0}.
\end{equation}
If
\[
  y_{1}(x_{0})y_{2}'(x_{0})-y_{1}'(x_{0})y_{2}(x_{0})=0,
\]
it's impossible to satisfy \eqref{eq:5.1.23} and \eqref{eq:5.1.24} (and therefore \eqref{eq:5.1.22})
unless $k_{0}$ and $k_{1}$ happen to satisfy
\begin{eqnarray*}
  y_{1}(x_{0})k_{1}-y_{1}'(x_{0})k_{0} & = & 0 \\
  y_{2}'(x_{0})k_{0}-y_{2}(x_{0})k_{1} & = & 0.
\end{eqnarray*}
On the other hand, if
\begin{equation}
  \label{eq:5.1.25}y_{1}(x_{0})y_{2}'(x_{0})-y_{1}'(x_{0})y_{2}(x_{0})\ne0
\end{equation}
we can divide \eqref{eq:5.1.23} and \eqref{eq:5.1.24} through by the quantity on the left to obtain
\begin{equation}
  \label{eq:5.1.26}
  \begin{array}{rcl}
    c_{1} & = & \dst\frac{y_{2}'(x_{0})k_{0}-y_{2}(x_{0})k_{1}}{y_{1}(x_{0})y_{2}'(x_{0})-y_{1}'(x_{0})y_{2}(x_{0})}  \\[4\jot]
    c_{2} & = & \dst\frac{y_{1}(x_{0})k_{1}-y_{1}'(x_{0})k_{0}}{y_{1}(x_{0})y_{2}'(x_{0})-y_{1}'(x_{0})y_{2}(x_{0})},
  \end{array}
\end{equation}
no matter how $k_{0}$ and $k_{1}$ are chosen. This motivates us to consider conditions on $y_{1}$
and $y_{2}$ that imply \eqref{eq:5.1.25}.

\begin{theorem}
  \color{blue}
  \label{thmtype:5.1.4} Suppose $p$ and $q$ are continuous on $(a,b)$, let $y_{1}$ and $y_{2}$ be
  solutions of
  \begin{equation}
    \label{eq:5.1.27}y''+p(x)y'+q(x)y=0
  \end{equation}
  on $(a,b)$, and define
  \begin{equation}
    \label{eq:5.1.28}W=y_{1}y_{2}'-y_{1}'y_{2}.
  \end{equation}
  Let $x_{0}$ be any point in $(a,b)$. Then
  \begin{equation}
    \label{eq:5.1.29}W(x)=W(x_{0})e^{-\int^{x}_{x_{0}}p(t)\, dt},\quad a<x<b.
  \end{equation}
  Therefore either $W$ has no zeros in $(a,b)$ or $W\equiv0$ on $(a,b)$.
\end{theorem}

\proof Differentiating \eqref{eq:5.1.28} yields
\begin{equation}
  \label{eq:5.1.30}W'=y'_{1}y'_{2}+y_{1}y''_{2}-y'_{1}y'_{2}-y''_{1}y_{2}= y_{1}y''_{2}-y''_{1}y_{2}.
\end{equation}
Since $y_{1}$ and $y_{2}$ both satisfy \eqref{eq:5.1.27},
\[
  y''_{1}=-py'_{1}-qy_{1}\mbox{\quad and \quad}y''_{2}=-py'_{2}-qy_{2}.
\]
Substituting these into \eqref{eq:5.1.30} yields
\begin{eqnarray*}
  W' & = & \dst-y_{1}\bigl(py'_{2}+qy_{2}\bigr)+y_{2}\bigl(py'_{1}+qy_{1}\bigr) \\
     & = & \dst-p(y_{1}y'_{2}-y_{2}y'_{1})-q(y_{1}y_{2}-y_{2}y_{1})             \\
     & = & -p(y_{1}y'_{2}-y_{2}y'_{1})=-pW.
\end{eqnarray*}
Therefore $W'+p(x)W=0$; that is, $W$ is the solution of the initial value problem
\[
  y'+p(x)y=0,\quad y(x_{0})=W(x_{0}).
\]
We leave it to you to verify by separation of variables that this implies \eqref{eq:5.1.29}. If
$W(x_{0})\ne0$, \eqref{eq:5.1.29} implies that $W$ has no zeros in $(a,b)$, since an exponential is
never zero. On the other hand, if $W(x_{0})=0$, \eqref{eq:5.1.29} implies that $W(x)=0$ for all $x$
in $(a,b)$.\bbox

The function $W$ defined in \eqref{eq:5.1.28} is the
\href{https://mathshistory.st-andrews.ac.uk/Biographies/Wronski/}{\color{blue}
\it Wronskian of $\{y_{1},y_{2}\}$}. Formula \eqref{eq:5.1.29} is
\href{https://mathshistory.st-andrews.ac.uk/Biographies/Abel/}{\color{blue}
\it Abel's formula}.

The Wronskian of $\{y_{1},y_{2}\}$ is usually written as the determinant
\[
  W=\left|
  \begin{array}{cc}
    y_{1}  & y_{2}  \\[2\jot]
    y'_{1} & y'_{2}
  \end{array}\right|.
\]
The expressions in \eqref{eq:5.1.26} for $c_{1}$ and $c_{2}$ can be written in terms of determinants
as
\[
  c_{1}=\frac{1}{W(x_{0})}\left|
  \begin{array}{cc}
    k_{0} & y_{2}(x_{0})  \\[2\jot]
    k_{1} & y'_{2}(x_{0})
  \end{array}\right| \mbox{\quad and \quad}c_{2}=\frac{1}{W(x_{0})}\left|
  \begin{array}{cc}
    y_{1}(x_{0})  & k_{0} \\[2\jot]
    y'_{1}(x_{0}) & k_{1}
  \end{array}\right|.
\]
If you've taken linear algebra you may recognize this as
\href{https://mathshistory.st-andrews.ac.uk/Biographies/Cramer/}{\color{blue}
\it Cramer's rule}.

\begin{example}
  \label{example:5.1.5} \rm Verify Abel's formula for the following differential equations and the
  corresponding solutions, from Examples~\ref{example:5.1.1}, \ref{example:5.1.2}, and
  \ref{example:5.1.3}:
  \begin{alist}
    \item % (a)
    $y''-y=0;\quad y_{1}=e^{x},\;y_{2}=e^{-x}$

    \item % (b)
    $y''+\omega^{2}y=0;\quad\quad y_{1}=\cos\omega x,\;y_{2}=\sin\omega x$

    \item % (c)
    $x^{2}y''+xy'-4y=0;\quad y_{1}=x^{2},\;y_{2}=1/x^{2}$
  \end{alist}
\end{example}

\solutionpart{a} Since $p\equiv0$, we can verify Abel's formula by showing that $W$ is constant,
which is true, since
\[
  W(x)=\left|
  \begin{array}{rr}
    e^{x} & e^{-x}  \\[2\jot]
    e^{x} & -e^{-x}
  \end{array}\right|=e^{x}(-e^{-x})-e^{x}e^{-x}=-2
\]
for all $x$.

\solutionpart{b} Again, since $p\equiv0$, we can verify Abel's formula by showing that $W$ is
constant, which is true, since
\begin{eqnarray*}
  W(x)&=&\dst{\left|\begin{array}{cc}\cos\omega x&\sin\omega x\\[2\jot] -\omega\sin\omega x&\omega\cos\omega x\end{array}\right|}\\[2\jot] &=&\cos\omega x(\omega\cos\omega x)-(-\omega\sin\omega x)\sin\omega x\\&=&\omega(\cos^{2}\omega x+\sin^{2}\omega x)=\omega
\end{eqnarray*}
for all $x$.

\solutionpart{c} Computing the Wronskian of $y_{1}=x^{2}$ and $y_{2}=1/x^{2}$ directly yields
\begin{equation}
  \label{eq:5.1.31}W=\left|
  \begin{array}{cc}
    x^{2} & 1/x^{2}  \\[2\jot]
    2x    & -2/x^{3}
  \end{array}\right|=x^{2}\left(-\frac{2}{x^{3}}\right)-2x\left(\frac{1}{x^{2}}\right)=-\frac{4}{x}.
\end{equation}
To verify Abel's formula we rewrite the differential equation as
\[
  y''+\frac{1}{x}y'-\frac{4}{x^{2}}y=0
\]
to see that $p(x)=1/x$. If $x_{0}$ and $x$ are either both in $(-\infty,0)$ or both in $(0,\infty)$
then
\[
  \int_{x_{0}}^{x}p(t)\,dt=\int_{x_{0}}^{x}\frac{dt}{t}=\ln\left(\frac{x}{x_{0}}\right),
\]
so Abel's formula becomes
\begin{eqnarray*}
  W(x) & = & W(x_{0})e^{-\ln(x/x_{0})}=W(x_{0})\frac{x_{0}}{x}                                            \\
       & = & -\left(\frac{4}{x_{0}}\right)\left(\frac{x_{0}}{x}\right)\mbox{\quad from \eqref{eq:5.1.31}} \\
       & = & -\frac{4}{x},
\end{eqnarray*}
which is consistent with \eqref{eq:5.1.31}.\bbox

The next theorem will enable us to complete the proof of Theorem~\ref{thmtype:5.1.3}.

\begin{theorem}
  \color{blue}
  \label{thmtype:5.1.5} Suppose $p$ and $q$ are continuous on an open interval $(a,b)$, let $y_{1}$
  and $y_{2}$ be solutions of
  \begin{equation}
    \label{eq:5.1.32}y''+p(x)y'+q(x)y=0
  \end{equation}
  on $(a,b)$, and let $W=y_{1}y_{2}'-y_{1}'y_{2}$. Then $y_{1}$ and $y_{2}$ are linearly independent
  on $(a,b)$ if and only if $W$ has no zeros on $(a,b)$.
\end{theorem}

\proof We first show that if $W(x_{0})=0$ for some $x_{0}$ in $(a,b)$, then $y_{1}$ and $y_{2}$ are
linearly dependent on $(a,b)$. Let $I$ be a subinterval of $(a,b)$ on which $y_{1}$ has no zeros.
(If there's no such subinterval, $y_{1}\equiv0$ on $(a,b)$, so $y_{1}$ and $y_{2}$ are linearly
independent, and we're finished with this part of the proof.) Then $y_{2}/y_{1}$ is defined on $I$,
and
\begin{equation}
  \label{eq:5.1.33}\left(\frac{y_{2}}{y_{1}}\right)'=\frac{y_{1}y_{2}'-y_{1}'y_{2}}{y_{1}^{2}}=\frac{W}{y_{1}^{2}}.
\end{equation}
However, if $W(x_{0})=0$, Theorem~\ref{thmtype:5.1.4} implies that $W\equiv0$ on $(a,b)$. Therefore
\eqref{eq:5.1.33} implies that $(y_{2}/y_{1})'\equiv0$, so $y_{2}/y_{1}=c$ (constant) on $I$. This
shows that $y_{2}(x)=cy_{1}(x)$ for all $x$ in $I$. However, we want to show that $y_{2}=cy_{1}(x)$
for all $x$ in $(a,b)$. Let $Y=y_{2}-cy_{1}$. Then $Y$ is a solution of \eqref{eq:5.1.32} on $(a,b)$
such that $Y\equiv0$ on $I$, and therefore $Y'\equiv0$ on $I$. Consequently, if $x_{0}$ is chosen
arbitrarily in $I$ then $Y$ is a solution of the initial value problem
\[
  y''+p(x)y'+q(x)y=0,\quad y(x_{0})=0,\quad y'(x_{0})=0,
\]
which implies that $Y\equiv0$ on $(a,b)$, by the paragraph following Theorem~\ref{thmtype:5.1.1}.
(See also Exercise~\ref{exer:5.1.24}). Hence, $y_{2}-cy_{1}\equiv0$ on $(a,b)$, which implies that
$y_{1}$ and $y_{2}$ are not linearly independent on $(a,b)$.

Now suppose $W$ has no zeros on $(a,b)$. Then $y_{1}$ can't be identically zero on $(a,b)$ (why
not?), and therefore there is a subinterval $I$ of $(a,b)$ on which $y_{1}$ has no zeros. Since
\eqref{eq:5.1.33} implies that $y_{2}/y_{1}$ is nonconstant on $I$, $y_{2}$ isn't a constant
multiple of $y_{1}$ on $(a,b)$. A similar argument shows that $y_{1}$ isn't a constant multiple of
$y_{2}$ on $(a,b)$, since
\[
  \left(\frac{y_{1}}{y_{2}}\right)'=\frac{y_{1}'y_{2}-y_{1}y_{2}'}{y_{2}^{2}}=-\frac{W}{y_{2}^{2}}
\]
on any subinterval of $(a,b)$ where $y_{2}$ has no zeros. \bbox

We can now complete the proof of Theorem~\ref{thmtype:5.1.3}. From Theorem~\ref{thmtype:5.1.5}, two
solutions $y_{1}$ and $y_{2}$ of \eqref{eq:5.1.32} are linearly independent on $(a,b)$ if and only
if $W$ has no zeros on $(a,b)$. From Theorem~\ref{thmtype:5.1.4} and the motivating comments
preceding it, $\{y_{1},y_{2}\}$ is a fundamental set of solutions of \eqref{eq:5.1.32} if and only
if $W$ has no zeros on $(a,b)$. Therefore $\{y_{1},y_{2}\}$ is a fundamental set for
\eqref{eq:5.1.32} on $(a,b)$ if and only if $\{y_{1},y_{2}\}$ is linearly independent on $(a,b)$.
\bbox

The next theorem summarizes the relationships among the concepts discussed in this section.

\begin{theorem}
  \color{blue}
  \label{thmtype:5.1.6} Suppose $p$ and $q$ are continuous on an open interval $(a,b)$ and let
  $y_{1}$ and $y_{2}$ be solutions of
  \begin{equation}
    \label{eq:5.1.34}y''+p(x)y'+q(x)y=0
  \end{equation}
  on $(a,b)$. Then the following statements are equivalent; that is, they are either all true or all
  false.
  \begin{alist}
    \item % (a)
    The general solution of \eqref{eq:5.1.34} on $(a,b)$ is $y=c_{1}y_{1}+c_{2}y_{2}$.

    \item % (b)
    $\{y_{1},y_{2}\}$ is a fundamental set of solutions of \eqref{eq:5.1.34} on $(a,b)$.

    \item % (c)
    $\{y_{1},y_{2}\}$ is linearly independent on $(a,b)$.

    \item % (d)
    The Wronskian of $\{y_{1},y_{2}\}$ is nonzero at some point in $(a,b)$.

    \item % (d)
    The Wronskian of $\{y_{1},y_{2}\}$ is nonzero at all points in $(a,b)$.
  \end{alist}
\end{theorem}

We can apply this theorem to an equation written as
\[
  P_{0}(x)y''+P_{1}(x)y'+P_{2}(x)y=0
\]
on an interval $(a,b)$ where $P_{0}$, $P_{1}$, and $P_{2}$ are continuous and $P_{0}$ has no zeros.

\begin{theorem}
  \color{blue}
  \label{thmtype:5.1.7} Suppose $c$ is in $(a,b)$ and $\alpha$ and $\beta$ are real numbers, not
  both zero. Under the assumptions of Theorem~\ref{thmtype:5.1.7}, suppose $y_{1}$ and $y_{2}$ are
  solutions of \eqref{eq:5.1.34} such that
  \begin{equation}
    \label{eq:5.1.35}\alpha y_{1}(c)+\beta y_{1}'(c)=0\text{\; and\; \; }\alpha y_{2}(c)+\beta y_{2}'(c)=0.
  \end{equation}
  Then $\{y_{1},y_{2}\}$ isn't linearly independent on $(a,b)$.
\end{theorem}

\proof Since $\alpha$ and $\beta$ are not both zero, \eqref{eq:5.1.35} implies that
\[
  \left|
  \begin{array}{ccccccc}
    y_{1}(c) & y_{1}'(c) \\
    y_{2}(c) & y_{2}'(c)
  \end{array}\right|=0,\text{\; so\; \; }\left|
  \begin{array}{cccccc}
    y_{1}(c)  & y_{2}(c)  \\
    y_{1}'(c) & y_{2}'(c)
  \end{array}\right|=0
\]
and Theorem~\ref{thmtype:5.1.6} implies the stated conclusion.

\exercises
\begin{exerciselist}
  \item\label{exer:5.1.1}
  \begin{alist}
    \item % (a)
    Verify that $y_{1}=e^{2x}$ and $y_{2}=e^{5x}$ are solutions of
    \[
      y''-7y'+10y=0 \eqno{\rm(A)}
    \]
    on $(-\infty,\infty)$.

    \item % (b)
    Verify that if $c_{1}$ and $c_{2}$ are arbitrary constants then $y=c_{1}e^{2x}+c_{2}e^{5x}$ is a
    solution of (A) on $(-\infty,\infty)$.

    \item % (c)
    Solve the initial value problem
    \[
      y''-7y'+10y=0,\quad y(0)=-1,\quad y'(0)=1.
    \]

    \item % (d)
    Solve the initial value problem
    \[
      y''-7y'+10y=0,\quad y(0)=k_{0},\quad y'(0)=k_{1}.
    \]
  \end{alist}

  \item\label{exer:5.1.2}
  \begin{alist}
    \item % (a)
    Verify that $y_{1}=e^{x}\cos x$ and $y_{2}=e^{x}\sin x$ are solutions of
    \[
      y''-2y'+2y=0 \eqno{\rm(A)}
    \]
    on $(-\infty,\infty)$.

    \item % (b)
    Verify that if $c_{1}$ and $c_{2}$ are arbitrary constants then
    $y=c_{1}e^{x}\cos x+c_{2}e^{x}\sin x$ is a solution of (A) on $(-\infty,\infty)$.

    \item % (c)
    Solve the initial value problem
    \[
      y''-2y'+2y=0,\quad y(0)=3,\quad y'(0)=-2.
    \]

    \item % (d)
    Solve the initial value problem
    \[
      y''-2y'+2y=0,\quad y(0)=k_{0},\quad y'(0)=k_{1}.
    \]
  \end{alist}

  \item\label{exer:5.1.3}
  \begin{alist}
    \item % (a)
    Verify that $y_{1}=e^{x}$ and $y_{2}=xe^{x}$ are solutions of
    \[
      y''-2y'+y=0 \eqno{\rm(A)}
    \]
    on $(-\infty,\infty)$.

    \item % (b)
    Verify that if $c_{1}$ and $c_{2}$ are arbitrary constants then $y=e^{x}(c_{1}+c_{2}x)$ is a
    solution of (A) on $(-\infty,\infty)$.

    \item % (c)
    Solve the initial value problem
    \[
      y''-2y'+y=0,\quad y(0)=7,\quad y'(0)=4.
    \]

    \item % (d)
    Solve the initial value problem
    \[
      y''-2y'+y=0,\quad y(0)=k_{0},\quad y'(0)=k_{1}.
    \]
  \end{alist}

  \item\label{exer:5.1.4}
  \begin{alist}
    \item % (a)
    Verify that $y_{1}=1/(x-1)$ and $y_{2}=1/(x+1)$ are solutions of
    \[
      (x^{2}-1)y''+4xy'+2y=0 \eqno{\rm(A)}
    \]
    on $(-\infty,-1)$, $(-1,1)$, and $(1,\infty)$. What is the general solution of (A) on each of
    these intervals?

    \item % (b)
    Solve the initial value problem
    \[
      (x^{2}-1)y''+4xy'+2y=0,\quad y(0)=-5,\quad y'(0)=1.
    \]
    What is the interval of validity of the solution?

    \item % (c)
    \CGex Graph the solution of the initial value problem.

    \item % (d)
    Verify Abel's formula for $y_{1}$ and $y_{2}$, with $x_{0}=0$.
  \end{alist}

  \item\label{exer:5.1.5} Compute the Wronskians of the given sets of functions.

  \begin{tabular}[t]{@{}p{168pt}@{}p{168pt}}
    {\bf (a)} $\{1,e^{x}\}$ & {\bf (b)} $\{e^{x},e^{x}\sin x\}$
  \end{tabular}

  \begin{tabular}[t]{@{}p{168pt}@{}p{168pt}}
    {\bf (c)} $\{x+1,x^{2}+2\}$ & {\bf (d)} $\{x^{1/2},x^{-1/3}\}$
  \end{tabular}

  \begin{tabular}[t]{@{}p{168pt}@{}p{168pt}}
    {\bf (e)} $\left\{\dst\frac{\sin x}{x},\frac{\cos x}{x}\right\}$ & {\bf (f)} $\{x\ln|x|,x^{2}\ln|x|\}$ \\
    {\bf (g)} $\{e^{x}\cos\sqrt{x},e^{x}\sin\sqrt{x}\}$              &
  \end{tabular}

  \item\label{exer:5.1.6} Find the Wronskian of a given set $\{y_{1},y_{2}\}$ of solutions of
  \[
    y''+3(x^{2}+1)y'-2y=0,
  \]
  given that $W(\pi)=0$.

  \item\label{exer:5.1.7} Find the Wronskian of a given set $\{y_{1},y_{2}\}$ of solutions of
  \[
    (1-x^{2})y''-2xy'+\alpha(\alpha+1)y=0,
  \]
  given that $W(0)=1$. (This is
  \href{https://mathshistory.st-andrews.ac.uk/Biographies/Legendre/}{\color{blue}
  \it Legendre's equation}.)

  \item\label{exer:5.1.8} Find the Wronskian of a given set $\{y_{1},y_{2}\}$ of solutions of
  \[
    x^{2}y''+xy'+(x^{2}-\nu^{2})y=0,
  \]
  given that $W(1)=1$. (This is
  \href{https://mathshistory.st-andrews.ac.uk/Biographies/Bessel/}{\color{blue}
  \it Bessel's equation}.)

  \item\label{exer:5.1.9} %\exerciseabel
  (This exercise shows that if you know one nontrivial solution of $y''+p(x)y'+q(x)y=0$, you can use
  Abel's formula to find another.)

  Suppose $p$ and $q$ are continuous and $y_{1}$ is a solution of
  \[
    y''+p(x)y'+q(x)y=0 \eqno{\rm(A)}
  \]
  that has no zeros on $(a,b)$. Let $P(x)=\int p(x)\,dx$ be any antiderivative of $p$ on $(a,b)$.
  \begin{alist}
    \item % (a)
    Show that if $K$ is an arbitrary nonzero constant and $y_{2}$ satisfies
    \[
      y_{1}y_{2}'-y_{1}'y_{2}=Ke^{-P(x)}\eqno{\rm(B)}
    \]
    on $(a,b)$, then $y_{2}$ also satisfies (A) on $(a,b)$, and $\{y_{1},y_{2}\}$ is a fundamental
    set of solutions on (A) on $(a,b)$.

    \item % (b)
    Conclude from \part{a} that if $y_{2}=uy_{1}$ where $u'=K\dst\frac{e^{-P(x)}}{y_{1}^{2}(x)}$,
    then $\{y_{1},y_{2}\}$ is a fundamental set of solutions of (A) on $(a,b)$.
  \end{alist}

  \exercisetext{In Exercises~\ref{exer:5.1.10}--\ref{exer:5.1.23} use the method suggested by
  Exercise~\ref{exer:5.1.9} to find a second solution $y_{2}$ that isn't a constant multiple of the
  solution $y_{1}$. Choose $K$ conveniently to simplify $y_{2}$.}

  \item\label{exer:5.1.10} $y''-2y'-3y=0$; \quad $y_{1}=e^{3x}$

  \item\label{exer:5.1.11} $y''-6y'+9y=0$; \quad $y_{1}=e^{3x}$

  \item\label{exer:5.1.12}%
  $y''-2ay'+a^{2}y=0$\; ($a=$ constant); \quad $y_{1}=e^{ax}$

  \item\label{exer:5.1.13} $x^{2}y''+xy'-y=0$; \quad $y_{1}=x$

  \item\label{exer:5.1.14} $x^{2}y''-xy'+y=0$; \quad $y_{1}=x$

  \item\label{exer:5.1.15} $x^{2}y''-(2a-1)xy'+a^{2}y=0$\; ($a=$ nonzero constant);\, $x>0$; \quad
        $y_{1}=x^{a}$

  \item\label{exer:5.1.16} $4x^{2}y''-4xy'+(3-16x^{2})y=0$; \quad $y_{1}=x^{1/2}e^{2x}$

  \item\label{exer:5.1.17} $(x-1)y''-xy'+y=0$; \quad $y_{1}=e^{x}$

  \item\label{exer:5.1.18} $x^{2}y''-2xy'+(x^{2}+2)y=0$; \quad $y_{1}=x\cos x$

  \item\label{exer:5.1.19} $4x^{2}(\sin x)y''-4x(x\cos x+\sin x)y'+(2x\cos x+3\sin x)y=0$; \quad
        $y_{1}=x^{1/2}$

  \item\label{exer:5.1.20} $(3x-1)y''-(3x+2)y'-(6x-8)y=0$; \quad $y_{1}=e^{2x}$

  \item\label{exer:5.1.21} $(x^{2}-4)y''+4xy'+2y=0$; \quad $y_{1}=\dst\frac{1}{x-2}$

  \item\label{exer:5.1.22} $(2x+1)xy''-2(2x^{2}-1)y'-4(x+1)y=0$;\quad $y_{1}=\dst\frac{1}{x}$

  \item\label{exer:5.1.23} $(x^{2}-2x)y''+(2-x^{2})y'+(2x-2)y=0$;\quad $y_{1}=e^{x}$

  \item\label{exer:5.1.24} Suppose $p$ and $q$ are continuous on an open interval $(a,b)$ and let
        $x_{0}$ be in $(a,b)$. Use Theorem~\ref{thmtype:5.1.1} to show that the only solution of the
        initial value problem
  \[
    y''+p(x)y'+q(x)y=0,\quad y(x_{0})=0,\quad y'(x_{0})=0
  \]
  on $(a,b)$ is the trivial solution $y\equiv0$.

  \item\label{exer:5.1.25} Suppose $P_{0}$, $P_{1}$, and $P_{2}$ are continuous on $(a,b)$ and let
        $x_{0}$ be in $(a,b)$. Show that if either of the following statements is true then
        $P_{0}(x)=0$ for some $x$ in $(a,b)$.
  \begin{alist}
    \item % (a)
    The initial value problem
    \[
      P_{0}(x)y''+P_{1}(x)y'+P_{2}(x)y=0,\quad y(x_{0})=k_{0},\quad y'(x_{0})=k_{1}
    \]
    has more than one solution on $(a,b)$.

    \item % (b)
    The initial value problem
    \[
      P_{0}(x)y''+P_{1}(x)y'+P_{2}(x)y=0,\quad y(x_{0})=0,\quad y'(x_{0})=0
    \]
    has a nontrivial solution on $(a,b)$.
  \end{alist}

  \item\label{exer:5.1.26} Suppose $p$ and $q$ are continuous on $(a,b)$ and $y_{1}$ and $y_{2}$ are
        solutions of
  \[
    y''+p(x)y'+q(x)y=0 \eqno{\rm(A)}
  \]
  on $(a,b)$. Let
  \[
    z_{1}=\alpha y_{1}+\beta y_{2}\mbox{\quad and \quad}z_{2}=\gamma y_{1}+\delta y_{2},
  \]
  where $\alpha$, $\beta$, $\gamma$, and $\delta$ are constants. Show that if $\{z_{1},z_{2}\}$ is a
  fundamental set of solutions of (A) on $(a,b)$ then so is $\{y_{1},y_{2}\}$.

  \item\label{exer:5.1.27} Suppose $p$ and $q$ are continuous on $(a,b)$ and $\{y_{1},y_{2}\}$ is a
        fundamental set of solutions of
  \[
    y''+p(x)y'+q(x)y=0 \eqno{\rm(A)}
  \]
  on $(a,b)$. Let
  \[
    z_{1}=\alpha y_{1}+\beta y_{2}\mbox{\quad and \quad}z_{2}=\gamma y_{1}+\delta y_{2},
  \]
  where $\alpha$, $\beta$, $\gamma$, and $\delta$ are constants. Show that $\{z_{1},z_{2}\}$ is a
  fundamental set of solutions of (A) on $(a,b)$ if and only if $\alpha\gamma-\beta\delta\ne0$.

  \item\label{exer:5.1.28} Suppose $y_{1}$ is differentiable on an interval $(a,b)$ and
        $y_{2}=ky_{1}$, where $k$ is a constant. Show that the Wronskian of $\{y_{1},y_{2}\}$ is
        identically zero on $(a,b)$.

  \item\label{exer:5.1.29} Let
  \[
    y_{1}=x^{3}\quad\mbox{ and }\quad y_{2}=\left\{
    \begin{array}{rl}
      x^{3},  & x\ge0, \\
      -x^{3}, & x<0.
    \end{array}\right.
  \]
  \begin{alist}
    \item % (a)
    Show that the Wronskian of $\{y_{1},y_{2}\}$ is defined and identically zero on
    $(-\infty,\infty)$.

    \item % (b)
    Suppose $a<0<b$. Show that $\{y_{1},y_{2}\}$ is linearly independent on $(a,b)$.

    \item % (b)
    Use Exercise~\ref{exer:5.1.25}\part{b} to show that these results don't contradict
    Theorem~\ref{thmtype:5.1.5}, because neither $y_{1}$ nor $y_{2}$ can be a solution of an
    equation
    \[
      y''+p(x)y'+q(x)y=0
    \]
    on $(a,b)$ if $p$ and $q$ are continuous on $(a,b)$.
  \end{alist}

  \item\label{exer:5.1.30} Suppose $p$ and $q$ are continuous on $(a,b)$ and $\{y_{1},y_{2}\}$ is a
        set of solutions of
  \[
    y''+p(x)y'+q(x)y=0
  \]
  on $(a,b)$ such that either $y_{1}(x_{0})=y_{2}(x_{0})=0$ or $y_{1}'(x_{0})=y_{2}'(x_{0})=0$ for
  some $x_{0}$ in $(a,b)$. Show that $\{y_{1},y_{2}\}$ is linearly dependent on $(a,b)$.

  \item\label{exer:5.1.31} Suppose $p$ and $q$ are continuous on $(a,b)$ and $\{y_{1},y_{2}\}$ is a
        fundamental set of solutions of
  \[
    y''+p(x)y'+q(x)y=0
  \]
  on $(a,b)$. Show that if $y_{1}(x_{1})=y_{1}(x_{2})=0$, where $a<x_{1}<x_{2}<b$, then $y_{2}(x)=0$
  for some $x$ in $(x_{1},x_{2})$. \hint{Show that if $y_{2}$ has no zeros in $(x_{1},x_{2})$, then
  $y_{1}/y_{2}$ is either strictly increasing or strictly decreasing on $(x_{1},x_{2})$, and deduce
  a contradiction.}

  \item\label{exer:5.1.32} Suppose $p$ and $q$ are continuous on $(a,b)$ and every solution of
  \[
    y''+p(x)y'+q(x)y=0 \eqno{\rm(A)}
  \]
  on $(a,b)$ can be written as a linear combination of the twice differentiable functions
  $\{y_{1},y_{2}\}$. Use Theorem~\ref{thmtype:5.1.1} to show that $y_{1}$ and $y_{2}$ are themselves
  solutions of (A) on $(a,b)$.

  \item\label{exer:5.1.33} Suppose $p_{1}$, $p_{2}$, $q_{1}$, and $q_{2}$ are continuous on $(a,b)$
        and the equations
  \[
    y''+p_{1}(x)y'+q_{1}(x)y=0\mbox{\quad and \quad}y''+p_{2}(x)y'+q_{2}(x)y=0
  \]
  have the same solutions on $(a,b)$. Show that $p_{1}=p_{2}$ and $q_{1}=q_{2}$ on $(a,b)$.
  \hint{Use Abel's formula.}

  \item\label{exer:5.1.34} (For this exercise you have to know about $3\times3$ determinants.) Show
        that if $y_{1}$ and $y_{2}$ are twice continuously differentiable on $(a,b)$ and the
        Wronskian $W$ of $\{y_{1},y_{2}\}$ has no zeros in $(a,b)$ then the equation
  \[
    \frac{1}{W}\left|
    \begin{array}{ccc}
      y   & y_{1}   & y_{2}   \\[2\jot]
      y'  & y'_{1}  & y'_{2}  \\[2\jot]
      y'' & y_{1}'' & y_{2}''
    \end{array}\right|=0
  \]
  can be written as
  \[
    y''+p(x)y'+q(x)y=0, \eqno{\rm(A)}
  \]
  where $p$ and $q$ are continuous on $(a,b)$ and $\{y_{1},y_{2}\}$ is a fundamental set of
  solutions of (A) on $(a,b)$. \hint{Expand the determinant by cofactors of its first column.}

  \item\label{exer:5.1.35} Use the method suggested by Exercise~\ref{exer:5.1.34} to find a linear
        homogeneous equation for which the given functions form a fundamental set of solutions on
        some interval.

  \begin{tabular}[t]{@{}p{168pt}@{}p{168pt}}
    {\bf (a)} $e^{x}\cos2x,\quad e^{x}\sin2x$ & {\bf (b)} $x,\quad e^{2x}$
  \end{tabular}

  \begin{tabular}[t]{@{}p{168pt}@{}p{168pt}}
    {\bf (c)} $x,\quad x\ln x$ & {\bf (d)} $\cos(\ln x),\quad\sin(\ln x)$
  \end{tabular}

  \begin{tabular}[t]{@{}p{168pt}@{}p{168pt}}
    {\bf (e)} $\cosh x,\quad\sinh x$ & {\bf (f)} $x^{2}-1,\quad x^{2}+1$
  \end{tabular}

  \item\label{exer:5.1.36} Suppose $p$ and $q$ are continuous on $(a,b)$ and $\{y_{1},y_{2}\}$ is a
        fundamental set of solutions of
  \[
    y''+p(x)y'+q(x)y=0 \eqno{\rm(A)}
  \]
  on $(a,b)$. Show that if $y$ is a solution of (A) on $(a,b)$, there's exactly one way to choose
  $c_{1}$ and $c_{2}$ so that $y=c_{1}y_{1}+c_{2}y_{2}$ on $(a,b)$.

  \item\label{exer:5.1.37} Suppose $p$ and $q$ are continuous on $(a,b)$ and $x_{0}$ is in $(a,b)$.
        Let $y_{1}$ and $y_{2}$ be the solutions of
  \[
    y''+p(x)y'+q(x)y=0 \eqno{\rm(A)}
  \]
  such that
  \[
    y_{1}(x_{0})=1,\quad y'_{1}(x_{0})=0\mbox{\quad and \quad}y_{2}(x_{0})=0,\;y'_{2}(x_{0})=1.
  \]
  (Theorem~\ref{thmtype:5.1.1} implies that each of these initial value problems has a unique
  solution on $(a,b)$.)
  \begin{alist}
    \item % (a)
    Show that $\{y_{1},y_{2}\}$ is linearly independent on $(a,b)$.

    \item % (b)
    Show that an arbitrary solution $y$ of (A) on $(a,b)$ can be written as
    $y=y(x_{0})y_{1}+y'(x_{0})y_{2}$.

    \item % (c)
    Express the solution of the initial value problem
    \[
      y''+p(x)y'+q(x)y=0,\quad y(x_{0})=k_{0},\quad y'(x_{0})=k_{1}
    \]
    as a linear combination of $y_{1}$ and $y_{2}$.
  \end{alist}

  \item\label{exer:5.1.38} Find solutions $y_{1}$ and $y_{2}$ of the equation $y''=0$ that satisfy
        the initial conditions
  \[
    y_{1}(x_{0})=1,\quad y'_{1}(x_{0})=0\mbox{\quad and \quad}y_{2}(x_{0})=0, \quad y'_{2}(x_{0})=1.
  \]
  Then use Exercise~\ref{exer:5.1.37}~{\bf (c)} to write the solution of the initial value problem
  \[
    y''=0,\quad y(0)=k_{0},\quad y'(0)=k_{1}
  \]
  as a linear combination of $y_{1}$ and $y_{2}$.

  \item\label{exer:5.1.39} Let $x_{0}$ be an arbitrary real number. Given
        (Example~\ref{example:5.1.1}) that $e^{x}$ and $e^{-x}$ are solutions of $y''-y=0$, find
        solutions $y_{1}$ and $y_{2}$ of $y''-y=0$ such that
  \[
    y_{1}(x_{0})=1,\quad y'_{1}(x_{0})=0\mbox{\quad and \quad}y_{2}(x_{0})=0,\;y'_{2}(x_{0})=1.
  \]
  Then use Exercise~\ref{exer:5.1.37}~{\bf (c)} to write the solution of the initial value problem
  \[
    y''-y=0,\quad y(x_{0})=k_{0},\quad y'(x_{0})=k_{1}
  \]
  as a linear combination of $y_{1}$ and $y_{2}$.

  \item\label{exer:5.1.40} Let $x_{0}$ be an arbitrary real number. Given
        (Example~\ref{example:5.1.2}) that $\cos\omega x$ and $\sin\omega x$ are solutions of
        $y''+\omega^{2}y=0$, find solutions of $y''+\omega^{2}y=0$ such that
  \[
    y_{1}(x_{0})=1,\quad y'_{1}(x_{0})=0\mbox{\quad and \quad}y_{2}(x_{0})=0,\;y'_{2}(x_{0})=1.
  \]
  Then use Exercise~\ref{exer:5.1.37}~{\bf (c)} to write the solution of the initial value problem
  \[
    y''+\omega^{2}y=0,\quad y(x_{0})=k_{0},\quad y'(x_{0})=k_{1}
  \]
  as a linear combination of $y_{1}$ and $y_{2}$. Use the identities
  \begin{eqnarray*}
    \cos(A+B) & = & \cos A\cos B-\sin A\sin B \\
    \sin(A+B) & = & \sin A\cos B+\cos A\sin B
  \end{eqnarray*}
  to simplify your expressions for $y_{1}$, $y_{2}$, and $y$.

  \item\label{exer:5.1.41} Recall from Exercise~\ref{exer:5.1.4} that $1/(x-1)$ and $1/(x+1)$ are
        solutions of
  \[
    (x^{2}-1)y''+4xy'+2y=0 \eqno{\rm(A)}
  \]
  on $(-1,1)$. Find solutions of (A) such that
  \[
    y_{1}(0)=1,\quad y'_{1}(0)=0\mbox{\quad and \quad}y_{2}(0)=0,\;y'_{2}(0)=1.
  \]
  Then use Exercise~\ref{exer:5.1.37}~{\bf (c)} to write the solution of initial value problem
  \[
    (x^{2}-1)y''+4xy'+2y=0,\quad y(0)=k_{0},\quad y'(0)=k_{1}
  \]
  as a linear combination of $y_{1}$ and $y_{2}$.

  \item\label{exer:5.1.42}
  \begin{alist}
    \item % (a)
    Verify that $y_{1}=x^{2}$ and $y_{2}=x^{3}$ satisfy
    \[
      x^{2}y''-4xy'+6y=0 \eqno{\rm(A)}
    \]
    on $(-\infty,\infty)$ and that $\{y_{1},y_{2}\}$ is a fundamental set of solutions of (A) on
    $(-\infty,0)$ and $(0,\infty)$.

    \item % (b)
    Let $a_{1}$, $a_{2}$, $b_{1}$, and $b_{2}$ be constants. Show that
    \[
      y=\left\{
      \begin{array}{rr}
        a_{1}x^{2}+a_{2}x^{3}, & x\ge0,         \\
        b_{1}x^{2}+b_{2}x^{3}, & x<0\phantom{,}
      \end{array}\right.
    \]
    is a solution of (A) on $(-\infty,\infty)$ if and only if $a_{1}=b_{1}$. From this, justify the
    statement that $y$ is a solution of (A) on $(-\infty,\infty)$ if and only if
    \[
      y=\left\{
      \begin{array}{rr}
        c_{1}x^{2}+c_{2}x^{3}, & x\ge0, \\
        c_{1}x^{2}+c_{3}x^{3}, & x<0,
      \end{array}\right.
    \]
    where $c_{1}$, $c_{2}$, and $c_{3}$ are arbitrary constants.

    \item % (c)
    For what values of $k_{0}$ and $k_{1}$ does the initial value problem
    \[
      x^{2}y''-4xy'+6y=0,\quad y(0)=k_{0},\quad y'(0)=k_{1}
    \]
    have a solution? What are the solutions?

    \item % (d)
    Show that if $x_{0}\ne0$ and $k_{0}$, $k_{1}$ are arbitrary constants, the initial value problem
    \[
      x^{2}y''-4xy'+6y=0,\quad y(x_{0})=k_{0},\quad y'(x_{0})=k_{1}\eqno{\rm(B)}
    \]
    has infinitely many solutions on $(-\infty,\infty)$. On what interval does (B) have a unique
    solution?
  \end{alist}

  \item\label{exer:5.1.43} %\exercisec
  \begin{alist}
    \item % (a)
    Verify that $y_{1}=x$ and $y_{2}=x^{2}$ satisfy
    \[
      x^{2}y''-2xy'+2y=0 \eqno{\rm(A)}
    \]
    on $(-\infty,\infty)$ and that $\{y_{1},y_{2}\}$ is a fundamental set of solutions of (A) on
    $(-\infty,0)$ and $(0,\infty)$.

    \item % (b)
    Let $a_{1}$, $a_{2}$, $b_{1}$, and $b_{2}$ be constants. Show that
    \[
      y=\left\{
      \begin{array}{rr}
        a_{1}x+a_{2}x^{2}, & x\ge0,         \\
        b_{1}x+b_{2}x^{2}, & x<0\phantom{,}
      \end{array}\right.
    \]
    is a solution of (A) on $(-\infty,\infty)$ if and only if $a_{1}=b_{1}$ and $a_{2}=b_{2}$. From
    this, justify the statement that the general solution of (A) on $(-\infty,\infty)$ is
    $y=c_{1}x+c_{2}x^{2}$, where $c_{1}$ and $c_{2}$ are arbitrary constants.

    \item % (c)
    For what values of $k_{0}$ and $k_{1}$ does the initial value problem
    \[
      x^{2}y''-2xy'+2y=0,\quad y(0)=k_{0},\quad y'(0)=k_{1}
    \]
    have a solution? What are the solutions?

    \item % (d)
    Show that if $x_{0}\ne0$ and $k_{0}$, $k_{1}$ are arbitrary constants then the initial value
    problem
    \[
      x^{2}y''-2xy'+2y=0,\quad y(x_{0})=k_{0},\quad y'(x_{0})=k_{1}
    \]
    has a unique solution on $(-\infty,\infty)$.
  \end{alist}

  \item\label{exer:5.1.44}
  \begin{alist}
    \item % (a)
    Verify that $y_{1}=x^{3}$ and $y_{2}=x^{4}$ satisfy
    \[
      x^{2}y''-6xy'+12y=0 \eqno{\rm(A)}
    \]
    on $(-\infty,\infty)$, and that $\{y_{1},y_{2}\}$ is a fundamental set of solutions of (A) on
    $(-\infty,0)$ and $(0,\infty)$.

    \item % (b)
    Show that $y$ is a solution of (A) on $(-\infty,\infty)$ if and only if
    \[
      y=\left\{
      \begin{array}{rr}
        a_{1}x^{3}+a_{2}x^{4}, & x\ge0, \\
        b_{1}x^{3}+b_{2}x^{4}, & x<0,
      \end{array}\right.
    \]
    where $a_{1}$, $a_{2}$, $b_{1}$, and $b_{2}$ are arbitrary constants.

    \item % (c)
    For what values of $k_{0}$ and $k_{1}$ does the initial value problem
    \[
      x^{2}y''-6xy'+12y=0,\quad y(0)=k_{0},\quad y'(0)=k_{1}
    \]
    have a solution? What are the solutions?

    \item % (d)
    Show that if $x_{0}\ne0$ and $k_{0}$, $k_{1}$ are arbitrary constants then the initial value
    problem
    \[
      x^{2}y''-6xy'+12y=0,\quad y(x_{0})=k_{0},\quad y'(x_{0})=k_{1}\eqno{\rm(B)}
    \]
    has infinitely many solutions on $(-\infty,\infty)$. On what interval does (B) have a unique
    solution?
  \end{alist}
\end{exerciselist}
